\subsection{Motivación y posición cronológica}

Checkpoint~03D extiende la línea global de Checkpoint~03, aunque se ejecutó
cronológicamente después de Checkpoint~05. La numeración se conserva porque el experimento
responde una pregunta pendiente de la primera entrega global: si los grids deliberadamente
pequeños de 03C.2 eran la causa principal de la limitación de los expertos globales.

Los 59 rendimientos FIRA permanecen completamente ocultos.

\subsection{Representaciones, familias y presupuesto}

Se evaluaron cuatro representaciones deterministas:
\[
G0=\text{base},\quad
G1=\text{agronomic},\quad
G2=\text{base+agronomic},\quad
G3=\text{all deterministic}.
\]
No se reabrió descubrimiento target-aware. El universo incluyó CatBoost y XGBoost en GPU,
LightGBM, ExtraTrees, HistGradientBoosting y controles Ridge/PLS, para 23 pares
representación--familia.

La primera configuración wide se detuvo después de siete outer fits de CatBoost completados,
con el octavo iniciado. Cada outer fit tomaba aproximadamente 15 minutos porque contenía
nueve candidatos, tres inner folds y un refit. Extrapolar ese costo a 230 outer fits implicaba
varios días de cómputo para 138 etiquetas.

La reducción afectó sólo tuning redundante. Se conservaron cinco folds outer state-stratified,
cinco municipality-grouped y tres inner folds. Cada familia no lineal pesada quedó con tres
candidatos representativos; boosting se limitó a 2,500 iteraciones/estimadores y ExtraTrees a
1,500 árboles. Los partials wide no se reutilizaron y \artifact{--resume} ahora verifica un
fingerprint del contrato del run.

\subsection{Resultados}

El run balanced terminó en 248.085 minutos, con 230 outer fits y 6,348 filas OOF.

\begin{table}[htbp]
\centering
\caption{Finalistas globales de Checkpoint 03D, ordenados por peor RMSE OOF entre los dos
protocolos históricos.}
\label{tab:cp03d-finalists}
\begin{tabular}{cllrr}
\toprule
Rank & Representación & Modelo & State RMSE & Grouped RMSE \\
\midrule
1 & G1 agronomic & HistGBLarge & 0.549924 & \textbf{0.715339} \\
2 & G1 agronomic & XGBoostLarge & \textbf{0.525847} & 0.733080 \\
3 & G1 agronomic & LightGBMLarge & 0.540812 & 0.748478 \\
\bottomrule
\end{tabular}
\end{table}

Los tres finalistas usan G1. Las concatenaciones más grandes no fueron más robustas: CatBoost
G2/G3 alcanzó state RMSE 0.5074/0.5146 pero grouped 0.9307/0.9305; ExtraTrees G2/G3 obtuvo
state 0.4988/0.5039 y grouped 0.8577/0.8812.

Los ensambles históricos E123 y E13 estaban alrededor de
$(0.5121,0.7481)$ y $(0.5082,0.7488)$ en (state, grouped). HistGB G1 mueve la frontera grouped
hasta 0.7153, pero empeora state. XGBoost G1 ofrece 0.5258/0.7331. Por tanto, el cómputo
adicional sí mueve parte de la frontera, pero no produce dominancia simultánea en ambos
protocolos.

\subsection{Deployment ONNX}

El ranking científico se separó del soporte de conversión. HistGB G1 es rank 1, pero el stack
estable no produjo un ONNX válido para esa familia. XGBoost G1, rank 2, es el finalista
deployable de mayor rango.

La verificación final sobre 138 filas obtuvo:
\[
\max_i |\hat y_i^{Python}-\hat y_i^{ONNX}|
=
3.814697\times10^{-6}.
\]
Los artefactos locales son:
\begin{center}
\artifact{models/final/checkpoint03d_global.joblib}\\
\artifact{models/final/checkpoint03d_global.onnx}\\
\artifact{models/final/checkpoint03d_global.onnx.json}.
\end{center}

\subsection{Consecuencia}

Los RMSE state/grouped de 03D y el RMSE target-matched de Local04D pertenecen a protocolos
distintos; no se comparan como si fueran una sola tabla. 03D cierra la búsqueda global amplia,
pero no sustituye el predictor competition-facing. Local04D permanece como método final
post-05. Cualquier mezcla con los nuevos globales debe evaluarse en un 05B explícito bajo las
mismas pseudo-competiciones target-matched.
